\section*{Bab \ref{ch:b2:alkene-redox} --- Oksidasi dan Reduksi Alkena}

\begin{solution}{exo:b2:alkene-redox:1}
Metilsikloheksana; 1-metilsikloheksan-1-ol; trans-2-metilsikloheksan-1-ol
(adisi sin H dan B, sehingga \ce{OH} dan H yang baru berada pada sisi yang
sama, dan metilnya trans terhadap \ce{OH}); 1-metilsikloheksena oksida;
cis-1-metilsikloheksana-1,2-diol; 6-oksoheptanal
\ce{CH3CO(CH2)4CHO}.
\end{solution}

\begin{solution}{exo:b2:alkene-redox:2}
Hidroborasi--oksidasi: 2-metilpropan-1-ol. Hidrasi terkatalisis asam:
2-metilpropan-2-ol.
\end{solution}

\begin{solution}{exo:b2:alkene-redox:3}
trans-2,3-Dimetiloksirana, dengan kedua metil pada sisi cincin yang
berseberangan. Senyawa ini kiral (tidak ada bidang cermin) dan terbentuk
sebagai campuran rasemat.
\end{solution}

\begin{solution}{exo:b2:alkene-redox:4}
Pengolahan reduktif: dua molekul propanal. Pengolahan oksidatif: dua molekul
asam propanoat.
\end{solution}

\begin{solution}{exo:b2:alkene-redox:5}
(a) meso-butana-2,3-diol, tidak aktif karena akiral. (b) diol-diol
(2\textit{R},3\textit{R}) dan (2\textit{S},3\textit{S}) dalam jumlah yang
sama, tidak aktif karena rasemat.
\end{solution}

\begin{solution}{exo:b2:alkene-redox:6}
Dalam asam, air menyerang karbon tersier; dalam basa, hidroksida menyerang
\ce{CH2}; dalam kedua kasus, satu \ce{OH} berakhir pada setiap karbon:
2-metilpropana-1,2-diol yang sama. Dengan metanol, kedua jalur itu
menghasilkan eter yang berbeda (seperti dalam bab ini).
\end{solution}

\begin{solution}{exo:b2:alkene-redox:7}
Hidrogenasi dengan satu ekuivalen \ce{H2} pada katalis paladium yang
diracuni (Lindlar): adisi sin pada ikatan rangkap tiga berhenti pada alkena
(\textit{Z}).
\end{solution}

\begin{solution}{exo:b2:alkene-redox:8}
Propanon (3 C) dan butanal (4 C) bergabung melalui karbon karbonilnya
masing-masing, sehingga terbentuk \ce{(CH3)2C=CH-CH2CH2CH3}, yaitu
2-metilheks-2-ena.
\end{solution}

\begin{solution}{exo:b2:alkene-redox:9}
\ce{BH3} (atau boran yang meruah), lalu \ce{H2O2} dan \ce{NaOH}: heksan-1-ol.
Hidrasi terkatalisis asam melalui karbokation sekunder dan menghasilkan
heksan-2-ol (Markovnikov), dengan produk-produk samping hasil penataan ulang.
\end{solution}

\begin{solution}{exo:b2:alkene-redox:10}
\omterm{def:b2:alkene-redox:epoxide}{Asam peroksi}, sebuah elektrofil, bereaksi lebih cepat dengan ikatan rangkap
cincin yang lebih tersubstitusi dan lebih kaya elektron. Boran yang meruah
bereaksi dengan \ce{C=CH2} ujung yang kurang terhalang.
\end{solution}

\begin{solution}{exo:b2:alkene-redox:11}
trans: \omterm{def:b2:alkene-redox:epoxide}{asam peroksi}, lalu air dalam asam (pembukaan anti). cis: \ce{OsO4}
katalitik dengan pengoksidasi ulang, atau permanganat encer dingin (sin).
\end{solution}

\begin{solution}{exo:b2:alkene-redox:12}
\ce{C6H10} mempunyai dua derajat ketidakjenuhan; satu ekuivalen \ce{H2}
berarti satu \ce{C=C}, jadi satu cincin. Satu dialdehida berkarbon enam
berasal dari cincin yang terbuka pada ikatan rangkapnya: sikloheksena.
\end{solution}

\begin{solution}{pb:b2:alkene-redox:1}
\textbf{1.} $10 \times 12.011 + 16 \times 1.008 = \qty{136.24}{g/mol}$.
\textbf{2.} $(2 \times 10 + 2 - 16)/2 = 3$.
\textbf{3.} Satu cincin dan dua ikatan rangkap.
\textbf{4.} \ce{C10H16 + 2H2 -> C10H20}, 1-isopropil-4-metilsikloheksana.
\textbf{5.} $1.00/136.24 = \qty{7.34}{mmol}$ limonena, \qty{14.68}{mmol} \ce{H2}.
\textbf{6.} $V = nRT/p = 0.01468 \times 8.314 \times 298.15/10^5 = \qty{3.64e-4}{m^3}$,
yaitu \qty{364}{mL}.
\textbf{7.} Tidak: karbon cincin 1 dan 4 masing-masing membawa dua jalur
cincin yang identik; isomer cis dan trans sama-sama akiral.
\textbf{8.} Metanal \ce{HCHO} (1 C) dari ujung \ce{=CH2}, dan satu senyawa
\ce{C9}, 3-asetil-6-oksoheptanal \ce{CH3CO-CH2CH2-CH(COCH3)-CH2-CHO}.
\textbf{9.} Kedua karbonnya termasuk dalam cincin yang sama: memutus ikatan
rangkap membuka cincin tanpa memisahkan apa pun.
\textbf{10.} Metanal (sebuah gas).
\textbf{11.} Gugus-gugus aldehida menjadi asam karboksilat (metanal berlanjut
menjadi \ce{CO2}); keton-ketonnya tidak berubah.
\textbf{12.} Keton dan aldehida rantai \ce{C9} dulu tersambung di dalam
cincin; keton kedua (\ce{COCH3}) dan metanal adalah kedua ujung ikatan
rangkap rantai samping.
\textbf{13.} Satu (C4 cincin); pusat itu tidak tersentuh oleh \omterm{def:b2:alkene-redox:cleavage}{ozonolisis} dan
bertahan sebagai pusat stereogenik produk \ce{C9}.
\textbf{14.} \ce{C=CH2} yang eksosiklik, yang kurang terhalang daripada ikatan
rangkap cincin yang trisubstitusi.
\textbf{15.} 2-(4-metilsikloheks-3-en-1-il)propan-1-ol: \ce{CH2OH} pada bekas
karbon ujung.
\textbf{16.} Primer.
\textbf{17.} \ce{OH} masuk ke karbon yang kurang tersubstitusi, berlawanan
dengan orientasi yang diramalkan aturan Markovnikov untuk hidrasi.
\textbf{18.} Dua konfigurasi pada pusat yang baru, yang menghasilkan dua
diastereomer (bersama pusat yang sudah ada).
\textbf{19.} Hidrasi terkatalisis asam (atau hidrasi Markovnikov lainnya):
alkohol tersier, $\alpha$-terpineol.
\textbf{20.} Ikatan rangkap cincin: trisubstitusi, lebih kaya elektron.
\textbf{21.} 1,2-Epoksi-1-metil-4-(prop-1-en-2-il)sikloheksana, dengan
oksigen yang menjembatani C1 dan C2.
\textbf{22.} Pada C1, karbon yang lebih tersubstitusi (yang membawa metil).
\textbf{23.} 1-Metil-4-(prop-1-en-2-il)sikloheksana-1,2-diol, dengan kedua
\ce{OH} trans (pembukaan anti).
\textbf{24.} \omterm{def:b2:alkene-redox:dihydroxylation}{Dihidroksilasi} sin pada ikatan rangkap cincin (\ce{OsO4}
katalitik), yang menghasilkan diol cis --- asalkan pereaksinya dapat dibuat
bereaksi secara selektif dengan ikatan rangkap itu.
\textbf{25.} \qty{1.00}{g} limonena menyerap $\boldsymbol{\approx \qty{364}{mL}}$
dihidrogen.
\end{solution}
